\section{Введение}

Огромная часть машинного обучения основана на решении задачи оптимизации без ограничений

\begin{equation}
    \label{eq:general}
	\min_{w \in \mathbb{R}^d} f(w).
\end{equation}


Задачи вида \eqref{eq:general} охватывают множество приложений, включая минимизацию эмпирического риска~\citep{chapelle2000vicinal},
глубокое обучение~\citep{lecun2015deep} и задачи обучения с учителем~\citep{cunningham2008supervised}, такие как метод наименьших квадратов с регуляризацией~\citep{rifkin2007notes} или логистическая регрессия~\citep{shalev2014understanding}.

Классический метод решения задачи оптимизации \eqref{eq:general} --- градиентный спуск:

\begin{equation}
    \label{GD}
	w_{t+1} = w_t - \eta \nabla f(w_t),
\end{equation}
где $\eta > 0$ --- размер шага (шаг обучения).

Задача минимизации \eqref{eq:general} может быть трудноразрешимой, особенно когда размер выборки или размерность задачи крайне велики.

В таких случаях подсчёт полного градиента на каждой итерации градиентного спуска становится очень дорогим по времени и вычислительным ресурсам, особенно учитывая, что градиентному спуску часто требуется большое количество итераций для сходимости.
В современном машинном обучении, особенно с появлением глубокого обучения, растёт интерес к решению всё более крупных и сложных задач. Популярным решением для таких проблем стал стохастический градиентный спуск~\citep{robbins1951stochastic, bottou2010large}.


С течением времени методы оптимизации постоянно совершенствовались и развивались, становясь всё более сложными.
Одним из важнейших аспектов в оптимизации является правильный подбор размера шага в ходе итерационного процесса.
Адаптивные методы с градиентным масштабированием динамически регулируют
этот размер шага на основе информации о градиенте. Такое адаптивное поведение, применяемое к каждой переменной, улучшает процесс оптимизации, позволяя эффективно перемещаться по сложным ландшафтам и обеспечивая оптимальный прогресс для каждой переменной~\citep{duchi2011adaptive}. В частности, эти методы приобрели значительную популярность в области машинного обучения, где преобладают высокоразмерные задачи~\citep{zhang2018three, yao2021adahessian}.

Более подробно, под методами с масштабированным градиентом понимаются техники, предполагающие предобуславливание градиента задачи некоторой матрицей $D_t$, что позволяет учитывать геометрию задачи. В общем случае шаг алгоритмов с предобуславливанием может быть выражен как следующая модификация шага \eqref{GD}:
\begin{equation}
    w_{t+1} = w_t - \eta \cdot D_t^{-1} g_t,
\end{equation}
где $g_t$ --- несмещённый стохастический градиент.

Идея использования матрицы масштабирования отсылает нас к методу Ньютона, где $D_t = \nabla^2 f(w_t)$, и квазиньютоновским методам~\citep{dennis1977quasi}. Однако вычисление и обращение гессиана сопряжено со значительными трудностями, что приводит нас к необходимости использования определённых эвристик в качестве замены матрицы $D_t$. Примерами таких эвристических методов являются Adagrad~\citep{duchi2011adaptive}, Adam~\citep{kingma2014adam}, RMSProp~\citep{rmsprop}, OASIS~\citep{jahani2021doubly} и так далее, где стратегии вычисления $D_t$ не требуют точной оценки гессиана.
Например, в Adagrad предобуславливатель представлен в виде:
$$ D_t = \text{diag} \left\{ \sqrt{\sum_{t'=0}^t{g_{t'} \odot g_{t'}}} \right\}, $$
где $\odot$ --- адамарово (покомпонентное) произведение. Отметим, что этот подход использует только стохастические градиенты.

RMSProp и Adam используют похожие идеи:
$$ D_t^2 = \beta D_{t-1}^2 + (1 - \beta) \text{diag} \left\{ g_{t} \odot g_{t} \right\}, $$
где $\beta \in (0, 1)$ представляет собой степень учёта предыдущих итераций~\citep{kingma2014adam}.

В OASIS используется другой подход:
$$ D_t = \text{diag} \left\{ z_t \odot \nabla^2 f(w_t) z_t \right\}, $$
где $z_t$ --- случайный вектор из распределения Радемахера, т.е. каждая компонента $z_t^i$ принимает значения $\pm 1$ с вероятностью $\frac{1}{2}$~\citep{jahani2021doubly}. На первый взгляд кажется, что требуется полный гессиан, однако в действительности он не вычисляется явно: произведение $\nabla^2 f(w_t) z_t$ находится дифференцированием скалярной функции $\langle \nabla f(w_t), z_t \rangle$, а $\text{diag}\{z_t \odot \nabla^2 f(w_t) z_t\}$ является несмещённой оценкой диагонали гессиана (оценка Хатчинсона).

Несмотря на преимущества методов предобуславливания, они склонны к переобучению~\citep{wilson2017marginal}; таким образом, возникает необходимость в их совместном применении с регуляризацией. Этот подход широко применяется для решения различных задач машинного обучения, включая классификацию изображений~\citep{zhu2017learning}, распознавание речи~\citep{zhou2017improved} и обработку естественного языка~\citep{wu2022stgn}, и показал свою эффективность в улучшении обобщающей способности нейронных сетей~\citep{girosi1995regularization}.

С регуляризацией задача \eqref{eq:general} переформулируется как
\begin{equation} \label{F_big}
	\min_{w \in \mathbb{R}^d} F(w) := f(w) + r(w),
\end{equation}
где $r$ --- функция регуляризации.

В методах с предобуславливанием есть несколько способов добавления регуляризации.
Можно включить регуляризатор $r$ в вычисление $g_t$, и тогда он будет учитываться при вычислении $D_t$. Этот способ равносилен применению базового метода к оптимизационной задаче \eqref{F_big}.
Также можно добавить регуляризатор непосредственно в шаг обновления весов, уменьшая норму $w$~\citep{loshchilov2017decoupled}.
Такой способ регуляризации называется затуханием весов и, что примечательно, оказывается более эффективным в практических задачах.
Существует и промежуточный способ учёта регуляризатора, который будет рассмотрен далее в работе.

Несмотря на свою практическую эффективность, методы, использующие затухание весов, относительно мало изучены с точки зрения теории сходимости.
В связи с этим возникает ряд исследовательских вопросов:
\begin{itemize}
    \item \textit{Сходятся ли с теоретической точки зрения методы с предобуславливанием и затуханием весов?}
    \item \textit{Если сходятся, то какова скорость их сходимости?}
    \item \textit{К какой задаче они сходятся?}
\end{itemize}
